\documentclass[../main.tex]{subfiles}
\begin{document}

\chapternotnumbered{Introduction}
\label{chap:Introduction}

The goal of this thesis is to investigate the use of certain numerical methods in simulating atmospheric flows around topography, commonly referred to as \emph{orographic flows}. The text is structured as follows:

\begin{compactenum}
    \item \Cref{chap:FLOWS} provides an introduction to the physical foundations and mathematical descriptions of orographic flows.
    \item \Cref{chap:MODELLING} presents the concept of well-balanced numerical schemes.
    \item \Cref{chap:SPH} summarizes the main numerical method used: Smoothed Particle Hydrodynamics (SPH). 
    \item \Cref{chap:EXPERIMENTS} presents numerical experiments obtained by the proposed methods.
\end{compactenum}

\section*{Main contributions}
While SPH has found its use in continuum mechanics and astrophysics, its application to stratified atmospheric flows remains underdeveloped. The methodological contributions of this thesis are:
\begin{compactenum}
    \item The derivation of a density-independent SPH formulation utilizing \emph{potential temperature} as the primary thermodynamic variable, aligning the method with the standard conventions of atmospheric sciences.
    \item The construction of a novel \emph{well-balanced} discretisation for this framework that preserves the hydrostatic background exactly at the discrete level. 
\end{compactenum}

\section*{Motivation}

Earth's atmosphere is one of the most thoroughly studied and well-described natural systems. With the development of computer architecture, a great effort has been made to complete the \emph{experimental} data with the \emph{predictive} abilities of computational modelling. The need for effective and reliable predictive power is naturally well-founded---many processes in the atmosphere have a massive influence on the lives of Earth's inhabitants.

One example of such a process is orographic flow---an atmospheric flow of air mass over topography. Such flows are fairly common in nature and most importantly they give rise to effects of high relevance in atmospheric dynamics \& weather forecasting, such as

\begin{compactenum}
	\item downslope windstorms;
	\item \textcolor{gray}{atmospheric} internal gravity waves\sidenotemark;
	\item clear air turbulence.
\end{compactenum}

\sidenotetext{Note the distinction between \emph{gravitational waves}---disturbances of \emph{space-time} studied in the general theory of relativity, and \emph{gravity waves}---disturbances of \emph{a fluid} studied mainly in meteorology and oceanology. To make this distinction even more evident, the name \emph{internal} gravity waves is sometimes used in the literature.}

The ability to obtain accurate and reliable \emph{numerical} solutions to such flows is thus highly \emph{practically relevant. \sidenotemark} \sidenotetext{And scientifically relevant also: the study of orographic flows and mainly internal gravity waves is an active field in the atmospheric science community.} In this work, we aim to explore the suitability of SPH for this task.


\end{document}
